\section{Phase 5: model and controller complexity}
\label{sec:phase5}

\subsection{Experiment 5A: LTI-prototype sweep}
\label{sec:experiment5a}
\paragraph{Objective and refinement of the original plan.}
After the independent-fleet validation in \cref{sec:experiment4g}, determine
how many frozen gains approximate M4's three scheduled families. The original
generic cluster-count sweep is refined to $n\in\{1,2,3,5,8,10\}$ speed anchors
per known family, hence $3n$ frozen controllers. This isolates speed
discretization while keeping structural specialization and identification fixed.
It is not a learned-clustering study.

\paragraph{Frozen setup and selection boundary.}
Reuse the ten 4G training/test seed pairs and their stored structured parameter
fits. Both M4 and LTI grids use these same fits and unchanged $Q,R$. For $n>1$,
candidate layouts are equally spaced endpoint-inclusive anchors over $[12,28]$
m/s, equal-cell centers, and training-speed quantiles at probabilities
$(j+1/2)/n$, $j=0,\ldots,n-1$. For $n=1$, compare 16, 20, and 24 m/s.
Each grid is shared across the three families. For every count, select the
amplitude-feasible candidate with smallest mean yaw RMSE on the two original
moderate training maneuvers and the training fleet. This is a finite
training-based search, not an optimal placement claim. No test metric enters
anchor selection.

Selected LTI grids use nearest-anchor switching, with the lower anchor at ties.
M4 retains interpolation over its five fixed anchors. Test the same three
scenarios and three amplitude factors as 4G on independent test clients.
Report tracking, steering RMS, peak steering rate, switch counts, amplitude
feasibility, and frozen stability. Neither method receives hysteresis,
smoothing, steering-rate limiting, or feedback-weight retuning.

\paragraph{What is counted.}
All methods originate from three fits with three parameter ratios each. Thus
this study cannot claim fewer independently learned models for M4. M4 deploys
15 gain vectors, each with three feedback coefficients and one prefilter,
plus five shared speed anchors: 65 scalars. A frozen grid deploys $3n$ vectors
and $n$ shared anchors: $13n$ scalars. Counts exclude common software, units,
labels, and metadata. Retaining explicit frozen $A,B$ matrices adds $18n$
matrix scalars. Alternatively, these matrices can be reconstructed from the
same nine identified ratios. Family count, gain-vector count, and storage
are therefore distinct quantities.

\paragraph{Matching rule.}
Moderate tracking averaged over all three scenarios is primary. A grid matches
M4 if the upper endpoint of a paired 95\% percentile-bootstrap interval for
$J_{\mathrm{LTI}}-1.05J_{\mathrm{M4}}$ is at most zero, all trajectories are
amplitude-feasible, and the worst test frozen radius is below one. Resample ten
complete seed pairs 10,000 times using seed 5101. The predeclared 5\% tolerance
does not imply exact equivalence. Steering-rate peaks are reported as a practical
cost without adding a new acceptance threshold. The minimum matching count is
descriptive over tested candidates, not a universal lower bound or a deployment
choice validated on a further untouched dataset.

\paragraph{Results, interpretation, and decision.}
The sweep completed 18,900 held-out closed-loop evaluations. Training selection
chooses the same layout in all ten seed pairs: 24 m/s for $n=1$, training
quantiles for $n=2,3,5$, and equal-cell centers for $n=8,10$.
\Cref{tab:experiment5a} reports moderate tracking and practical deployment costs.
All selected grids and M4 are amplitude-feasible at all three severities and
pass the held-out frozen audit, with worst radii below 0.967.

\begin{table}[htbp]
\centering
\caption{Experiment 5A: moderate results. RMSE is in degrees/s, and peak rate
is the maximum steering rate in degrees/s across held-out clients and scenarios.
Scalars include feedback gains, prefilters, and shared speed anchors.}
\label{tab:experiment5a}
\begin{tabular}{lrrrrr}
\toprule
Design & Gain vectors & Scalars & Mean RMSE & Gap to M4 & Peak rate\\
\midrule
M4 interpolation & 15 & 65 & 0.02757 & --- & 2.29\\
$n=1$ & 3 & 13 & 0.03181 & 15.4\% & 2.27\\
$n=2$ & 6 & 26 & 0.03310 & 20.1\% & 34.37\\
$n=3$ & 9 & 39 & 0.03271 & 18.7\% & 30.89\\
$n=5$ & 15 & 65 & 0.03393 & 23.1\% & 39.72\\
$n=8$ & 24 & 104 & 0.02925 & 6.1\% & 39.45\\
$n=10$ & 30 & 130 & 0.02834 & 2.8\% & 14.33\\
\bottomrule
\end{tabular}
\end{table}

The first tested matching grid has ten anchors per family, or 30 frozen gain
vectors. Its paired 95\% interval for the tolerance-adjusted tracking gap is
$[-0.000663,-0.000543]$ degrees/s, entirely below zero. Eight anchors per
family do not match: their interval is $[0.000273,0.000329]$ degrees/s.
This supports a twofold gain-storage benefit for the implemented interpolated
controller, not a tenfold storage benefit inferred from counting three families
against 30 frozen gains. Both methods use three structural families and the
same nine fitted parameter ratios.

\Cref{fig:experiment5a} shows that performance does not improve monotonically
with anchor count. The grids are not nested, the selected layouts differ, and
switches occur at different phases of a maneuver. Training-optimal selection
within a finite candidate list does not imply test-optimal placement. In
particular, the selected five-anchor grid is not the endpoint-inclusive M5
used in 4G. This prevents attributing the observed curve to model count alone.

\begin{figure}[htbp]
\centering
\includegraphics[width=\linewidth]{experiment_5a_complexity.pdf}
\caption{Experiment 5A: tracking versus frozen gain count and controller
storage, followed by peak steering rate. Error bars are seed standard deviations.
The horizontal tolerance line represents 5\% above M4's mean tracking error.}
\label{fig:experiment5a}
\end{figure}

The ten-anchor grid's steering RMS is 0.80576 degrees, nearly equal to M4's
0.80578 degrees, while peak steering rate rises from 2.29 to 14.33 degrees/s.
It makes up to 12 switches per varying-speed maneuver. At strong excitation,
the corresponding peak rates are 3.48 and 21.58 degrees/s. Comparable RMS
effort therefore does not imply comparable input smoothness.

Matching is aggregate, not universal across scenarios. With ten anchors per
family, gaps are 5.03\% on the original S-curve, 3.16\% on the shifted lane
change, and $-1.11\%$ at constant 22 m/s. The original S-curve narrowly fails
the scenario-specific 5\% interval criterion. A single 24 m/s gain already
slightly outperforms M4 on the constant-speed case. These findings reinforce
that the value of interpolation depends on the task and operating variation.

\paragraph{Decision and limitations.}
Within this frozen comparison, many more than three frozen gains are needed
to meet the aggregate tracking target, and M4 achieves smoother steering with
less controller storage than the first matching grid. The scheduling argument
is supported for the tested nearest-switching implementation. It is not a lower
bound over all LTI designs, optimized partitions, hysteretic switching, or
bumpless transfer schemes, and it proves no formal time-varying stability claim.
Phase 6 should now test where the advantage persists as structural heterogeneity
and speed range vary, before drawing a broader model-complexity conclusion.
